%!TEX root = ./../../../thesis-trier.tex

Here, we show the ideal intersection theorem (\cref{constr_sync:gen_results:thm:ideal_separation}).
Later, we mostly use this result
to deduce \PSPACE-hardness from known \PSPACE-hard problems, especially
for constraint languages written as unions (see Remark~\ref{constrsync:compl:ideal_intersection_for_unions} for details).

\begin{theorem}[Ideal Intersection Theorem]
	\label{constr_sync:gen_results:thm:ideal_separation}
	Let $I \subseteq \Sigma^*$ be a fixed ideal language that is regular.
	Suppose $L \subseteq \Sigma^*$ is any regular language,
	then
	$$
	\constrsync{L \cap I} \leACZEROm \constrsync{L}.
	$$
\end{theorem}
\begin{proof}
	By \cref{prelim:thm:sync_words_is_ideal}
	there exists an automaton $\mathcal A_I$
	such that $I = \Syn(\mathcal A_I)$.
	Let $\semiautomat{A}{Q}{\Sigma}{\delta}$ be an input to $\constrsync{L\cap I}$.
	By~\cref{prelim:eqn:product_aut_sync_words}, 
	it is $\Syn(\mathcal A) \cap I = \Syn(\mathcal A \times \mathcal A_I)$,
	which gives	
	\begin{equation}\label{constrsync:eqn:ideal-int}
	\Syn(\mathcal A) \cap (L \cap I) \ne \emptyset
	\Leftrightarrow 
    \Syn(\mathcal A \times \mathcal A_I) \cap L \ne \emptyset,
	\end{equation}
	i.e., $\mathcal A \times \mathcal A_I$ has a synchronizing
	word in $L$ if and only if $\mathcal A$ has a synchronizing word in $L \cap I$.
	As $I$ is fixed,
	we can construct the product automaton $\mathcal A \times \mathcal A_I$
	%for a fixed automaton $\mathcal A_I$ 
	by a circuit family in $\AC^0$ 
	 as we only have to add transitions
	and states to~$\mathcal A$ determined by a fixed automaton $\mathcal A_I$.
	So, by~\cref{constrsync:eqn:ideal-int} 
	constructing $\mathcal A \times \mathcal A_I$ gives a valid reduction to \constrsync{L}.	
\end{proof}

\begin{example}
	Let $\Sigma = \{a,b,c\}$ and consider the ideal language $I = \Sigma^* c^{n} \Sigma^*$
	that is regular.
	Let $L_1 = \{ u \in \Sigma^* : |u|_c \le n \}$,
	$L_2 = \{ u \in \Sigma^* : |u|_c = n \}$,
	and $L_3 = \{ u \in \Sigma^* : m \le |u|_c \le n \}$
	for $n,m \ge 1$ with $m \le n$.
	Then we have
	$L_1 \cap I = L_2 \cap I = L_3 \cap I = \{a,b\}^* c^n \{a,b\}^*$.
	From \crefitem{constr_sync:thm:introductory_examples}{constrsync:ex:item:abcL} 
	and \cref{constr_sync:gen_results:thm:ideal_separation} 
	we get that $\constrsync{L_i}$
	for $i\in \{1,2,3\}$
	are all \PSPACE-complete problems.
\end{example}

\begin{corollary}
	Let $L \subseteq \Sigma^*$ be regular
	and $u_1, \ldots, u_n \in \Sigma$.
	Then we have 
	\begin{align*}
	\constrsync{L\cap \Sigma^*u_1\cdots u_n\Sigma^*} & \leACZEROm \constrsync{L}, \\
	\constrsync{L\cap \Sigma^*u_1 \Sigma^* \cdots \Sigma^* u_n \Sigma^*} & \leACZEROm \constrsync{L}.
	\end{align*}
\end{corollary}
\begin{proof}
	The languages
	 $\Sigma^*u_1\cdots u_n\Sigma^*$
	 and $\Sigma^*u_1 \Sigma^* \cdots \Sigma^* u_n \Sigma^*$
	 are regular and ideals.
	 So, the claim follows from \cref{constr_sync:gen_results:thm:ideal_separation}.		
\end{proof}

\IN{principal ideal}
Mostly, we apply Theorem~\ref{constr_sync:gen_results:thm:ideal_separation}
with \emph{principal ideals} $\Sigma^* u \Sigma^*$
for $u \in \Sigma^*$.

As pointed out in the next remark, for a union the problem either has the same complexity as each part, or its complexity is bounded from below by the intersection of one part with a principal ideal.

\begin{remark}\label{constrsync:compl:ideal_intersection_for_unions}
	Consider two regular languages $L_1, L_2 \subseteq \Sigma^*$.
	There exists $u \in L_2$
	such that $u \notin \fact{L_1}$
	or $L_2 \subseteq \fact{L_1}$.
	By \cref{constrsync:genresults:cor:subset},
	in the former case, 
	we have
	\[
	\constrsync{L_1} \eqACZEROm \constrsync{L_1\cup L_2}
	\]
	and in the latter case, since then
	\[
	( L_1 \cup L_2 )  \cap \Sigma^* u \Sigma^* = L_2 \cap \Sigma^* u \Sigma^*,
	\]
	from \cref{constr_sync:gen_results:thm:ideal_separation}
	we get 
	$\constrsync{L_2 \cap \Sigma^*u \Sigma^*} \leACZEROm \constrsync{L_1\cup L_2}$.
	%	This scheme is often use in hardness proofs.
	%	However, finding an appropriate word $u \in \fact{L_2}$ 
	%	such that a provable 
	%	hard problem $\constrsync{L_2 \cap \Sigma^*u \Sigma^*}$
	%	results is the tricky part.
\end{remark}

Because we often use \cref{constr_sync:gen_results:thm:ideal_separation} to ``separate
a hard part out'' of a constraint language we called it the Ideal Separation Theorem
in~\cite{HoffmannCocoonExtended}, but we changed the name here because it involves an intersection
and we thought the new name to be more appropriate.

Lastly, we show how the Ideal Intersection Theorem  can be of help
in a heuristic way by employing it to exclude problematic
cases from a construction. This usage might not always be possible (in fact, in most cases it is not 
possible),
but we give the following example to also build some intuition
on various ways in which the Ideal Intersection Theorem 
can be applied (more examples of its application 
will follow in the rest of this work).

\begin{remark} \label{constrsync:rem:construction-ideal-intersection}
	Consider $L = \{bb,ba\}^*a\Gamma^*$ for $|\Gamma|\le 1$.
	We use a reduction from \SIS
	to show \PSPACE-hardness of $\constrsync{L}$.
	However, the construction we give
	does not work straightforward way for~$L$,
	but we then argue that it can be
	used for a modified problem obtained by intersection with an appropriate ideal.
	We then obtain \PSPACE-hardness of $\constrsync{L}$
	from \cref{constr_sync:gen_results:thm:ideal_separation}.
	
	Let $\semiautomat{A}{Q}{\{a,b\}}{\delta}$ and $S \subseteq Q$.
	Set $\semiautomat{A'}{Q \times \{\varepsilon,b\} \cup \{t\}}{\{a,b\}}{\delta'}$
	where~$\delta'(t, x) = t$ and
	\[
	\delta'((q, u), x) =
	\begin{cases}
		t                           & \text{if } u = \varepsilon \text{ and } x = a, \\
		(q, b) & \text{if } u = \varepsilon \text{ and } x = b, \\
		(\delta(q, a), \varepsilon) & \text{if } ux = bb, \\
		(\delta(q, b), \varepsilon) & \text{if } ux = ba, \\
		(q, \varepsilon)            & \text{if } ux \notin \{bb,ba\} \text{ and } q \notin S, \\
		t							& \text{if } ux \notin \{bb,ba\} \text{ and } q \in S, 		
	\end{cases}
	\]
	for $u \in \{\varepsilon, b\}$, $q \in Q$ and $x \in \{a,b\}$.
	
	First, suppose $u_1 \cdots u_n a v$
	with $u_1, \ldots, u_n \in \{bb,ba\}$ and $v \in \Gamma^*$ 
	synchronizes~$\mathcal A'$.
	Then it must map all states to~$t$.
	Consider the state $\delta'((q, \varepsilon), u_1\cdots u_n)$
	for some $q \in Q$.
	Then it must have the form $(p, \varepsilon)$
	for some $p \in Q$.
	Assume $p \notin S$.
	Then $\delta'((p, \varepsilon), a) = (p, \varepsilon)$
	and since $|\Gamma|\le 1$
	the word~$v$ cannot map $(p, \varepsilon)$ to~$t$.
	Thus we must have $p \in S$
	and then the word $v_1 \cdots v_n$
	with $v_i = a$ if $u_i = bb$
	and $v_i = b$ if $u_i = ba$
	maps~$Q$ into~$S$.	
	
	If $u_1 \cdots u_n$ for $u_1,\ldots,u_n \in \{a,b\}$
	maps~$Q$ into~$S$,
	then $v_1 \dots v_n a$ with $v_i = bb$ if $u_i = a$
	and $v_i = ba$ if $u_i = b$ for all $i \in \{1,\ldots,n\}$
	maps all states of~$Q \times \{\varepsilon\}$ to~$t$.
	However, if $u_1 = \cdots = u_n = bb$,
	then $\delta'(Q \times \{b\}, u_1\cdots u_n) \subseteq Q \times \{b\}$
	and thus $\delta'(Q \times \{b\}, u_1 \cdots u_n a) \cap \delta(Q, a) \times \{\varepsilon\} \ne \emptyset$.
	Now no word from~$\Gamma^*$ can map this state to~$t$.
	But this problem does not arise if we read~$ba$
	at least once, since this word maps all states in $Q \times \{b\}$
	to~$t$.
	Now, note that $L \cap \Sigma^* baa \Sigma^* = \{bb,ba\}baa\Gamma^*$.
	Thus, by \cref{constr_sync:gen_results:thm:ideal_separation} we obtain \PSPACE-hardness for $\constrsync{L}$
	if we can show \PSPACE-hardness
	for $\constrsync{\{bb,ba\}baa\Gamma^*}$
	and to establish \PSPACE-hardness of this problem
	we can use the above construction since the problematic case discussed above
	does not arise for $\{bb,ba\}baa\Gamma^*$.
\end{remark}